\documentclass[11pt,a4paper]{article}
\usepackage[utf8]{inputenc}
\usepackage[margin=2cm]{geometry}
\usepackage{amssymb}
\usepackage{tabularx}
\usepackage{array}
\usepackage{tcolorbox}
\tcbuselibrary{breakable}

\tcbset{
    sharp corners,
    colback=white,
    colframe=black,
    boxrule=0.8pt,
    left=6pt,
    right=6pt,
    top=6pt,
    bottom=6pt,
    breakable
}

\setlength{\parindent}{0pt}
\setlength{\parskip}{5pt}

\begin{document}

% سربرگ
\begin{center}
    {\large \textbf{Bachelor's Graduation Project Questionnaire}}\\[3pt]
    {\large \textbf{Department of Industrial Engineering and Management Systems}}
\end{center}

\vspace{0.4cm}

\noindent
\begin{tabularx}{\textwidth}{|l>{\raggedright\arraybackslash}X|l>{\raggedright\arraybackslash}X|}
\hline
\textbf{Supervisor:} Dr. Abbas Ahmadi & \textbf{Student Name:} Reza Eskandari \\ \hline
\textbf{Student ID:} 40125003 & \textbf{First Registration Semester:} 7th Semester \\ \hline
\textbf{Completed Credits:} 137 & \textbf{GPA:} 18.10 / 20 \\ \hline
\multicolumn{2}{|l|}{\textbf{Project Title:}} \\
\multicolumn{2}{|p{0.96\textwidth}|}{\textbf{Dynamic Pricing Under Competition Using Reinforcement Learning and Bayesian Hyperparameter Optimization}} \\ \hline
\multicolumn{2}{|p{0.96\textwidth}|}{\textbf{Project Type:} \quad Theoretical $\square$ \quad Applied $\square$ \quad Other: Theoretical-Applied $\blacksquare$} \\ \hline
\end{tabularx}

\vspace{-1pt}

\begin{tcolorbox}
\textbf{\large Project Summary}

\vspace{0.2cm}
\textbf{\large Introduction}

In recent years, the expansion of e-commerce and the intense competition among online sellers have made dynamic pricing one of the critical tools for revenue management and maintaining market share. In such environments, customers exhibit high sensitivity to price changes by making real-time price comparisons across different platforms, which forces sellers to continuously adjust their prices according to market conditions. Traditional pricing methods, which are mainly based on static models, fixed demand assumptions, or predefined rules, are not capable of effectively responding to these rapid fluctuations and increasing competition. Therefore, the need for flexible, self-learning approaches that can adapt to market dynamics has become increasingly evident.

In this context, Reinforcement Learning (RL), as an approach based on interaction with the environment, has attracted considerable attention in the field of dynamic pricing. The experience-based learning nature of RL allows it to gradually learn optimal pricing policies by observing the outcomes of past decisions, without requiring predefined demand models. Early studies have shown that even basic algorithms such as Q-Learning can discover complex relationships between price and demand and outperform traditional methods.

\vspace{0.3cm}
\textbf{\large Problem Statement}

Despite the growing importance of dynamic pricing in competitive markets, a substantial portion of existing research still relies on static models and simple pricing rules that are unable to respond effectively to rapid market changes and competitor behavior. Traditional methods often depend on fixed demand assumptions and cannot properly capture the complex interactions among price, demand, competition, and market dynamics. Consequently, there is a need for systems that can learn from experience and direct interaction with the environment and provide adaptive, intelligent, and efficient pricing policies. Studies such as \textit{Dynamic Retail Pricing via Q-Learning} have shown that even simple Reinforcement Learning algorithms can outperform classical models; however, these studies assume a single-agent market environment without the presence of competitors.

On the other hand, more recent studies such as \textit{Kastius \& Schlosser (2022)} demonstrate that introducing Reinforcement Learning into a competitive environment among sellers leads to much more complex behaviors. In duopolistic or multi-seller markets, learning agents must learn not only the relationship between price and demand but also the reactions of competitors. This research demonstrates that RL in a competitive environment can generate strategies that are substantially different from those in a single-agent environment and, in some cases, may lead to collusive-like patterns even without direct communication among sellers. This highlights the high sensitivity of RL performance to hyperparameters and structural settings. However, despite addressing this issue, the study does not provide any systematic mechanism for optimizing these hyperparameters.

Along the same lines, the study \textit{Leveraging Reinforcement Learning and Bayesian Optimization} demonstrates that Bayesian Optimization (BO) can intelligently tune the hyperparameters of RL algorithms and improve learning stability, convergence speed, and overall performance. Despite the success of this approach, the study was conducted in a non-competitive environment, and the impact of BO under conditions where competitor behavior and market dynamics play a decisive role was not investigated.

This research aims to address these gaps by combining the strengths of all three research streams and proposing a framework that can learn intelligent pricing policies using Reinforcement Learning in a simulated competitive environment while simultaneously leveraging Bayesian Optimization to improve the quality, stability, and efficiency of the learning process. Bridging this gap can lead to an effective tool for analyzing competitive behavior, increasing the revenues of online sellers, and gaining a deeper understanding of pricing dynamics in digital markets.

Therefore, the main novelty of this research lies in integrating three major dimensions---dynamic pricing, Reinforcement Learning, and Bayesian Optimization---into a comprehensive and unified framework. Although each of these areas has been studied independently in previous research, to the best of our knowledge, no study has simultaneously and coherently combined all three components in a competitive environment. The distinguishing feature of this research is the use of Bayesian Optimization to intelligently tune the hyperparameters of the RL algorithm under competitive conditions. While the application of BO in previous studies has primarily been limited to non-competitive environments, this research investigates, for the first time, its impact on the quality and stability of the learning process in the presence of competitors.

The characteristics of the competitive environment further reinforce the need for intelligent approaches. In such an environment, competitors' behavior directly affects the profit and reward of the learning agent, causing the market dynamics to become non-stationary. As a result, even minor changes in the hyperparameters of the Reinforcement Learning algorithm can have a significant impact on learning stability and quality. Studies such as \textit{Kastius \& Schlosser (2022)} have shown that competitive environments may give rise to complex behaviors such as aggressive undercutting, nonlinear responses, and even the formation of collusive-like patterns. Therefore, precise hyperparameter tuning is more critical in such environments than in non-competitive settings, and the use of approaches such as Bayesian Optimization can play a decisive role in improving the performance of the learning agent.

\vspace{0.3cm}
\textbf{\large Research Methodology}

The research methodology of this study is based on designing and analyzing a simulated competitive pricing environment and developing an RL framework optimized using Bayesian methods. The research process is designed in four main stages.

\textbf{Stage 1: Designing the Competitive Environment}\\
In the first stage, a competitive market consisting of one learning seller (the RL agent) and one or more competitors with predefined behaviors is simulated. This simulation is conducted to create a controlled and reproducible environment, allowing multiple controlled and repeatable experiments to be performed during the RL training process and enabling the effects of the environment design to be studied while ensuring the stability of the results.

\textbf{Stage 2: Developing the Reinforcement Learning Model}\\
In the second stage, a Reinforcement Learning algorithm is implemented to learn the optimal pricing policy. These algorithms are used to model optimal decision-making in a discrete environment. The main components of this stage include defining the State, Action, and Reward. This stage enables the agent to learn an optimal pricing policy in the presence of competitors through experience and trial-and-error. Initial studies such as \textit{Dynamic Retail Pricing via Q-Learning} have also shown that basic RL algorithms can outperform traditional methods even in single-agent environments.

\textbf{Stage 3: Bayesian Optimization of RL Hyperparameters}\\
In the third stage, given the sensitivity of RL performance to hyperparameters, Bayesian Optimization (BO) is employed to optimally tune them. The study \textit{Leveraging Reinforcement Learning and Bayesian Optimization} also demonstrates that BO can intelligently optimize the selection of hyperparameters for RL algorithms, resulting in improved stability, convergence speed, and overall model performance.

\textbf{Stage 4: Performance Evaluation}\\
In the final stage, the performance of the RL model optimized using BO is evaluated under different competitive conditions. This stage enables a detailed assessment of the effectiveness of the proposed framework and demonstrates how BO improves the quality of RL learning in competitive environments. The study by \textit{Kastius \& Schlosser (2022)} also emphasizes that introducing RL into competitive settings can lead to complex behaviors such as price attacks and even collusive-like patterns, further highlighting the importance of stability and appropriate hyperparameter tuning.

\vspace{0.3cm}
\textbf{Keywords:}\\
Dynamic Pricing, Reinforcement Learning, Bayesian Optimization, Competitive Market Simulation, Hyperparameter Tuning, Revenue Management, Single-Agent Reinforcement Learning, Demand Modeling
\end{tcolorbox}

\end{document}